\section{Additional Materials}

\subsection{Project Repositories}
\label{appendix:project-repositories}

The software packages produced for this thesis are available in the following repositories:

\begin{itemize}
    \item \textbf{Cuttlefish} for setup of Ceph clusters\\ \url{https://github.com/beanfacts/cuttlefish}
    \item \textbf{Scalie} for automated benchmarking of file system scaling performance with \acs{IOR} and mdtest\\ \url{https://github.com/beanfacts/scalie}
    \item \textbf{Ceph-Thesis} for the LaTeX source of this thesis, as well as supporting materials\\ \url{https://github.com/beanfacts/ceph-thesis}
\end{itemize}

\subsection{In-Memory tmpfs}
\label{appendix:tmpfs-in-memory}

In~\cref{sec:opa-bad-root-cause-analysis}, it was mentioned that tmpfs is an in-memory file system. This is true on the \ac{HPC} cluster nodes because the nodes are not configured with any swap space. However, tmpfs data may sometimes be swapped to disk if available. From the Linux manual page for tmpfs\footnote{\url{https://man7.org/linux/man-pages/man5/tmpfs.5.html}. Accessed 07.01.2026. This manual page can also be reproduced on most Linux systems by typing \texttt{man tmpfs} in the terminal.}, it is mentioned that \enquote{The filesystem can employ swap space when physical memory pressure demands it.}.

\subsection{Unprivileged Kernel Log Viewing}
\label{appendix:dmesg-unprivileged}

Whether \texttt{dmesg} is a privileged command depends on the \texttt{kernel.dmesg\_restrict} sysctl value. On certain Linux distributions and on the \ac{GWDG} \ac{HPC} cluster, this is set to 0 (unrestricted). As such, no privileges are required to run \texttt{dmesg} itself; however, configuration of the location of crash dumps remains restricted, necessitating the approach taken in~\cref{sec:opa-bad-root-cause-analysis}.

\pagebreak
\subsection{\acs{IOR}/mdtest Container Dockerfile}
\label{appendix:ior-mdtest-dockerfile}

\Cref{lst:ior-dockerfile} shows the Dockerfile used to build the container for running the \acs{IOR} and \texttt{mdtest} benchmarks. The container is based on the Ceph base image, which is currently based on CentOS Stream 9, and the required dependencies are installed and utilized to build the \acs{IOR} and \texttt{mdtest} binaries with support for \ac{RADOS} and \ac{CephFS}.

\lstinputlisting[
    caption={
        [Dockerfile for \acs{IOR}/mdtest Container]
        Dockerfile used to build the container for running the \acs{IOR} and \texttt{mdtest} benchmarks.
    },
    label={lst:ior-dockerfile},
    frame=single,
    breaklines=true
]{assets/cuttlefish/Dockerfile.ior}

\pagebreak
\subsection{Sample Ceph Keyring}
\label{appendix:sample-ceph-keyring}

\lstinputlisting[
    language=ini,
    caption={
        [Ceph Keyring File]
        A generated Ceph keyring file containing keys which are distributed 
        among deployed daemons in the cluster. Enough keys are generated in order
        to allow for the scaling up of the cluster without the need to regenerate 
        the keyring.
    },
    label={lst:cephkeyring},
    frame=single,
    breaklines=true
]{assets/cuttlefish/example.ceph.keyring}